\documentclass[12pt, oneside]{book}

% Encoding and Font Packages
\usepackage[T1]{fontenc} % Enables more efficient encoding for characters
\usepackage[utf8]{inputenc} % Allows for direct input of Unicode characters
\usepackage{microtype} % Improves typography by fine-tuning text spacing

% Font Setup for Official Text
\usepackage{charter} % Primary font for readability
\usepackage{bm} 
\usepackage{mathpazo} 

% Monospaced Font for Code Listings
\usepackage[scaled=.95]{inconsolata}

% Color and Graphics
\usepackage{array}
\usepackage{xcolor} % For custom color definitions
\usepackage{graphicx} % To insert images, diagrams, etc.
\usepackage{caption} % Customization of figure captions
\usepackage{subcaption} % For sub-figures (helpful for detailed diagrams)
\usepackage{tikz} % For high-quality, custom diagrams and illustrations
\usetikzlibrary{positioning, shapes.geometric, arrows, calc} % Libraries for complex diagram creation

% Margins and Page Layout
\usepackage[a4paper, margin=1in]{geometry} % Adjust page margins
\usepackage{setspace} % For line spacing
\setstretch{1.3} % Sets a nice line spacing for the text (readable and clear)

% Headers and Footers
\usepackage{fancyhdr} % For fancy headers and footers
\pagestyle{fancy}
\fancyhf{}
\fancyhead[L]{\nouppercase{\leftmark}} % Left-aligned chapter heading
\fancyhead[R]{\thepage} % Right-aligned page number
\fancypagestyle{plain}{
	\fancyhf{}
	\fancyfoot[C]{\thepage} % Centered page number for "plain" pages
}

% Hyperlink and Navigation
\usepackage{hyperref} % For hyperlinks and navigation within the document
\hypersetup{
	colorlinks=true, % Enables colored links
	linkcolor=black, % Links are in black (good for printing)
	citecolor=blue, % Citation links are blue
	urlcolor=cyan, % URLs are cyan
	pdfauthor={Mahdi}, % Author metadata for PDF
	pdftitle={Arliz}, % Title metadata for PDF
	pdfsubject={Programming, Arrays, and Data Structures}, % Subject for PDF
	pdfkeywords={Arrays, Data Structures, Programming, History of Computing}, % Keywords for PDF
}

% Bibliography
\usepackage[backend=biber,style=apa]{biblatex} % APA-style citations
\addbibresource{references.bib} % Add the bibliography file

% Listings for Code
\usepackage{listings} % For code listings and syntax highlighting
\lstset{
	basicstyle=\ttfamily\small, % Monospaced font for code
	frame=single, % Draws a box around the code block
	breaklines=true, % Automatically breaks lines in code
	numbers=left, % Line numbers on the left
	numberstyle=\tiny\color{gray}, % Style of line numbers
	keywordstyle=\color{myblue}\bfseries, % Keywords in blue and bold
	commentstyle=\color{olive}, % Comments in green
	stringstyle=\color{orange}, % Strings in orange
	backgroundcolor=\color{lightgray!20}, % Light background for code blocks
	captionpos=b, % Position the caption at the bottom
	escapeinside={(*@}{@*)}, % Allows for inline LaTeX inside code blocks
	morekeywords={array, structure, algorithm, complexity}, % Add relevant keywords for highlighting
}

% Table of Contents and Index
\usepackage{tocbibind} % Adds bibliography and index to the table of contents
\usepackage{imakeidx} % For index generation
\makeindex % Generates the index

% Title Customization
\usepackage{titlesec} % For customizing chapter and section styles
\titleformat{\chapter}[display]
{\normalfont\huge\bfseries}{\chaptername\ \thechapter}{20pt}{\Huge}
\titleformat{\section}{\Large\bfseries}{\thesection}{1em}{}
\titleformat{\subsection}{\large\bfseries}{\thesubsection}{1em}{}

% Algorithms and Pseudocode
\usepackage{algorithm} % For displaying algorithms
\usepackage{algpseudocode} % For pseudocode formatting
% Special boxes and environments
\usepackage{tcolorbox}
\tcbuselibrary{most}

\newtcolorbox{notebox}[1][]{
	colback=blue!5!white,
	colframe=blue!75!black,
	title=Note,
	#1
}

\newtcolorbox{tipbox}[1][]{
	colback=green!5!white,
	colframe=green!75!black,
	title=Tip,
	#1
}

\newtcolorbox{warningbox}[1][]{
	colback=red!5!white,
	colframe=red!75!black,
	title=Important,
	#1
}
% Multicolumn for Glossary or Index
\usepackage{multicol} % For multicolumn layouts (useful for glossary)

% Custom Colors
\definecolor{myblue}{RGB}{0, 102, 204} % Define a custom blue color for code and links
\definecolor{lightgray}{RGB}{240, 240, 240} % Light gray for code backgrounds

% Title and author
% \title{{\Huge\textbf{Arliz}}\\[0.5em]
% {\LARGE A Journey Through Arrays and Computer Fundamentals}\\[1em]
%	{\large From Bits to Data Structures}}
% \author{{\LARGE Mahdi}} % Author name (Large font)
% \date{{\large \today}} % Date of the document
\usepackage{pdfpages}
\begin{document}
	\includepdf[pages=-, pagecommand={}, width=\paperwidth, height=\paperheight]{logo.pdf}
	\frontmatter
	% Title Page (Front Matter)

\begin{titlepage}
	\begin{center}
		\vspace*{2cm}
		
		{\Huge\bfseries In Praise of Arliz}\\[2em]
		
		{\Large\scshape Mahdi}\\[0.5em]
		\textsc{Computer Engineering Student}\\[2em]
		
		\begin{minipage}{1\textwidth}
			\centering
			This book evolves. Every insight gained—whether a circuit, a structure,\\
				or a simple idea—is absorbed and integrated. Arliz is never complete.\\
				Because understanding never is.
		\end{minipage}
		
		\vfill
		
	
	{\large \textsc{First Edition}}\\[0.5em]
	{\large \today}
	
	\vspace*{1cm}
	{\small
		\textcopyright\ 2025 Mahdi “Genix”  
		\par
		Released under the MIT License  
		\par
		\url{https://github.com/m-mdy-m/Arliz}
	}
	\end{center}
\end{titlepage}

%----------------------------------------------
% Dedication
%----------------------------------------------
\cleardoublepage
\thispagestyle{empty}
\vspace*{6cm}
\begin{center}
	\emph{
		To those who build from first principles.\\
		To the silent thinkers who design before they speak.\\
		To the ones who see in systems—\\
		not just machines, but metaphors.\\
		This is for you.}
\end{center}
% Acknowledgments
\chapter*{Acknowledgments}
\thispagestyle{empty}
\addcontentsline{toc}{chapter}{Acknowledgments}
I would like to express my gratitude to everyone who supported me during the creation of this book. Special thanks to the open-source community for their invaluable resources and to all those who reviewed early drafts and provided feedback.

   %\maketitle
	\newpage
	\chapter*{Preface}
	\thispagestyle{empty}
	\addcontentsline{toc}{chapter}{Preface}
	
	Every book has its own story, and this book is no exception. If I were to summarize the process of creating this book in one word, that word would be “improvised.” Yet the truth is that Arliz is the result of pure, persistent curiosity that has grown in my mind for years. What you are reading now could be called a technical book, a collection of personal notes, or even a journal of unanswered questions and curiosities. But I—officially—call it a \emph{book}, because it is written not only for others but for myself, as a record of my learning journey and an effort to understand more precisely the concepts that once seemed obscure and, at times, frustrating.\\
	The story of Arliz began with a simple feeling: \textbf{curiosity}.  
	Curiosity about what an array truly is. Perhaps for many this question seems trivial, but for me this word—encountered again and again in algorithm and data structure discussions—always raised a persistent question.\\
	Every time I saw terms like \texttt{array}, \texttt{stack}, \texttt{queue}, \texttt{linked list}, \texttt{hash table}, or \texttt{heap}, I not only felt confused but sensed that something fundamental was missing. It was as if a key piece of the puzzle had been left out. The first brief, straightforward explanations I found in various sources never sufficed; they assumed you already knew exactly what an array is and why you should use it. But I was looking for the \emph{roots}. I wanted to understand from zero what an array means, how it was born, and what hidden capacities it holds.\\
	That realization led me to decide:  
	\emph{If I truly want to understand, I must start from zero.}\\	
	There is no deeper story behind the name “Arliz.” There is no hidden philosophy or special inspiration—just a random choice. I simply declared:  
	\emph{This book is called Arliz.}  
	You may pronounce it "Ar-liz," "Array-Liz," or any way you like. I personally say "ar-liz." That is all—simple and arbitrary.\\	
	But Arliz is not merely a technical book on data structures. In fact, \textbf{Arliz grows alongside me}. \\
	Whenever I learn something I deem worth writing, I add it to this book. Whenever I feel a section could be explained better or more precisely, I revise it. Whenever a new idea strikes me—an algorithm, an exercise, or even a simple diagram to clarify a structure—I incorporate it into Arliz.\\
	This means Arliz is a living project. As long as I keep learning, Arliz will remain alive.\\	
	The structure of this book has evolved around a simple belief: true understanding begins with context. That’s why Arliz doesn’t start with code or syntax, but with the origins of computation itself. We begin with the earliest tools and ideas—counting stones, the abacus, mechanical gears, and early notions of logic—long before transistors or binary digits came into play. From there, we follow the evolution of computing: from ancient methods of calculation to vacuum tubes and silicon chips, from Babbage’s Analytical Engine to the modern microprocessor. Along this journey, we discover that concepts like arrays aren’t recent inventions—they are the culmination of centuries of thought about how to structure, store, and process information.\\
	In writing this book, I have always tried to follow three principles:
	
	\begin{itemize}
		\item \textbf{Simplicity of Expression:} I strive to present concepts in the simplest form possible, so they are accessible to beginners and not superficial or tedious for experienced readers.
		\item \textbf{Concept Visualization:} I use diagrams, figures, and visual examples to explain ideas that are hard to imagine, because I believe visual understanding has great staying power.
		\item \textbf{Clear Code and Pseudocode:} Nearly every topic is accompanied by code that can be easily translated into major languages like C\texttt{++}, Java, or C\#, aiming for both clarity and practicality.
	\end{itemize}
	
	An important note: many of the algorithms in Arliz are implemented by myself. I did not copy them from elsewhere, nor are they necessarily the most optimized versions. My goal has been to understand and build them from scratch rather than memorize ready-made solutions. Therefore, some may run slower than standard implementations—or sometimes even faster. For me, the process of understanding and constructing has been more important than simply reaching the fastest result.\\	
	Finally, let me tell you a bit about myself:  
	I am \textbf{Mehdi}. If you prefer, you can call me by my alias: \emph{Genix}. I am a student of Computer Engineering (at least at the time of writing this). I grew up with computers—from simple games to typing commands in the terminal—and I have always wondered what lies behind this screen of black and green text. There is not much you need to know about me, just that I am someone who works with computers, sometimes gives them commands, and sometimes learns from them.\\	
	I hope this book will be useful for understanding concepts, beginning your learning journey, or diving deeper into data structures. \\	
	Arliz is freely available. You can access the PDF, LaTeX source, and related code at:  
	\begin{center}
		\url{https://github.com/m-mdy-m/Arliz}
	\end{center}
	In each chapter, I have included exercises and projects to aid your understanding. Please do not move on until you have completed these exercises, because true learning happens only by solving problems.\\	
	I hope this book serves you well—whether for starting out, reviewing, or simply satisfying your curiosity. And if you learn something, find an error, or have a suggestion, please let me know. As I said:
	\emph{This book grows with me.}
	\tableofcontents
	\renewcommand{\arraystretch}{1.5} % Adjust row height for better readability
	% Main Content
	\mainmatter

\part{The Birth of Computing: From Mechanical to Electronic}
\section*{Introduction}

Long before a single line of code was ever written—long before electricity, transistors, or even the concept of modern logic circuits—humans felt an innate drive to calculate, record, and model the world around them. Computing is not a recent invention. It is one of humanity’s oldest intellectual pursuits, rooted in necessity and evolved through creativity. Before we dive into complex abstractions like arrays or data structures, we must ask a deeper, almost philosophical question: \textbf{What does it mean to compute?}\\
This part of the book invites you on a journey—not just through the machinery and breakthroughs that brought us the modern computer, but through the evolution of human thought about numbers, representation, and control. Arrays, as we will later explore in depth, are not merely structures to store data. They are reflections of how we’ve ordered information for thousands of years. Their logic is built upon ancient insights—on sets, sequences, and patterns—and they embody the fundamental human need to represent, repeat, and manipulate structured information.\\
Our journey begins in ancient times, long before Christ, with devices like the abacus, first appearing over 2,500 years ago in Mesopotamia and later refined by Chinese, Roman, and Japanese cultures. The abacus was not just a calculator—it was an embodiment of the concepts of \textbf{state}, \textbf{position}, and \textbf{transformation}, principles that continue to underpin all modern computation. It allowed people to model quantities, track multiple values in parallel (an early echo of array indexing), and perform operations based on positional representation.\\
From these early tools, we progress into the classical mathematical age, where the Greeks formalized logic, and concepts like \textbf{sets} and \textbf{ordered lists} began to take philosophical shape. While not arrays in the modern sense, these ideas laid the intellectual groundwork for thinking about groups of data—grouped, related, or sequential—that could be acted upon as a whole. The set, in particular, became a foundational concept in mathematics and later in programming: an abstract container for elements that obey rules and enable operations. The leap from abstract sets to concrete arrays reflects one of the key transitions in computational history—from idea to implementation.\\
In the 17th century, visionaries like Blaise Pascal and Gottfried Wilhelm Leibniz attempted to automate arithmetic with mechanical devices. These weren’t just clever tools—they were the first signs of a dream to make thinking itself mechanical. Charles Babbage expanded this dream with his Analytical Engine in the 19th century, envisioning a machine that could be programmed and reprogrammed—a concept that wouldn’t become reality until a century later. Ada Lovelace, who worked with Babbage, went even further. She grasped that machines could go beyond numbers: they could process symbolic logic, follow instructions, and even imitate aspects of reasoning. She anticipated the algorithm as a mental construct, not just a set of steps.\\
As we move forward into the 20th century, the invention of electromechanical and electronic machines—using relays, vacuum tubes, and later transistors—marked a revolution. No longer limited by gears and levers, computers became faster, more reliable, and more abstract. The idea of a \textbf{stored program} emerged, allowing machines to modify their behavior dynamically. This wasn’t just a technical innovation—it was a conceptual transformation. Programs became data, and data became active. Arrays, now implemented in memory, could be changed, traversed, and manipulated at runtime—opening the door to software as we know it today.\\
Eventually, we arrive at logic gates, boolean algebra, and the transistor—the atomic units of modern computation. These are more than circuits; they are the physical embodiment of logical thought: conditions, branching, repetition. From gates we build circuits, from circuits microprocessors, and from those, machines that can simulate anything we can formalize.\\
Before concluding this part, we will look closely at how data is represented: binary numbers, encoding schemes, floating-point formats, and character representations. These are not just technical tools; they are perspectives. They define the limits of what a machine can know, express, and manipulate. And finally, we arrive at memory—where arrays live, grow, and function. Memory is not just storage; it is the canvas of computation. It is where change happens and where order emerges.\\
If you are excited to write code, build systems, and jump into implementation, you are free to skip ahead. But if you stay with us for this brief but essential historical and conceptual journey, you will see programming not just as control over a machine, but as part of a much older story: the story of how humans learned to structure thought, encode logic, and make abstract ideas come alive.\\
Let us begin—at the beginning. With sand, stone, wood, and brass. And with minds bold enough to imagine machines that think.
	\chapter{Mechanical Roots of Computing}
	\section{From the Abacus to the Analytical Engine}
	
	\subsection{The Abacus: The First Data Structure}
	
	\subsection{Pascalin and Leibniz's Wheel}
	
	\subsection{Babbie's Analytical Engine}
	
	\subsection{Ada Lovelace and the First Algorithm}
	
	\section{Electromechanical Computers and Early Concepts}
	
	\chapter{Introduction to Computers and Data Storage}
	\section{A Brief History of Computing}
	
	\subsection{From the Abacus to the Analytical Engine}
	
	\subsection{The Electronic Computer Revolution}
	
	\subsection{The Birth of Stored Programs}
	
	\chapter{The Birth of the Modern Computer and Its Architecture}
	\section{The Transition to Electronic Computing}
	
	\subsection{The Age of the Vacuum Tube}
	
	\subsection{ENIAC and Early Electronic Computers}
	
	\subsection{Von Neumann Architecture}
	
	\subsection{The Concept of a Program Saved}
	
	\chapter{Hardware Foundations}
	\section{Hardware Fundamentals}
	
	\subsection{Logic Circuits and Gates}
	
	\subsection{Von Neumann Architecture}
	
	\section{Logic Gates and Boolean Algebra}
	
	\section{Transistors: Building Blocks}
	
	\section{Integrated Circuits and Microprocessors}
	
	\section{Evolution of Computer Architecture}
	\section{The Birth of Modern Computer Architecture}
	
	\chapter{Digital Logic and Boolean Foundations}
	\section{Transistors: The Atomic Units of Computation}
	\section{Logic Gates and Circuit Design}
	\section{From NAND to NOR: Building Computational Primitives}
	
	\chapter{Number Systems and Data Representation}
	\section{Historic Counting Systems}
	\section{Binary: The Language of Machines}
	\subsection{Unsigned Integer Representation}
	\subsection{Two's Complement System}
	\section{Floating Point: Representing the Continuous}
	\section{Character Encoding Evolution}
	\subsection{From EBCDIC to Unicode}
	
	\chapter{Memory: The Computer's Canvas}
	\section{Historic Storage Media}
	\subsection{Punch Cards to Core Memory}
	\section{Modern Memory Hierarchy}
	\subsection{Registers and Cache Architecture}
	\subsection{RAM Geometries and Bank Organization}
	\section{Address Space Concepts}
	\subsection{Physical vs. Virtual Addressing}
	\subsection{Memory Mapping and Address Translation}
	\section{Memory Layout in Programs}
	\subsection{Code Segment (Text Segment)}
	\subsection{Data Segment (Initialized Data)}
	\subsection{BSS Segment (Uninitialized Data)}
	\subsection{Heap and Stack Segments}
	\section{Memory Allocation Strategies}
	\subsection{Static Memory Allocation}
	\subsection{Stack-Based Allocation}
	\subsection{Heap-Based Dynamic Allocation}
	\subsection{Memory Pools and Custom Allocators}
	\section{Memory Protection and Segmentation}
	\subsection{Memory Protection Mechanisms}
	\subsection{Segmentation and Paging}
	\subsection{Memory Management Unit (MMU)}
	\subsection{Address Space Layout Randomization (ASLR)}
		
	\part{The Array Odyssey}
	 \chapter{Historical Emergence of Arrays}
	 \section{Early Array Concepts in Mathematics}
	 \section{Arrays in Assembly Language}
	 \subsection{IBM 704 Index Registers}
	 \section{Array Adoption in High-Level Languages}
	 
	 \chapter{Array Anatomy}
	 \section{Formal Mathematical Definition}
	 \section{Machine Representation}
	 \subsection{Contiguous Memory Layout}
	 \subsection{Stride and Cache Considerations}
	 \section{Dimensionality Perspectives}
	 \subsection{Physical vs. Logical Dimensions}
	 
	 \chapter{Memory Layout Engineering}
	 \section{Static Allocation Strategies}
	 \subsection{BSS vs. DATA Segments}
	 \section{Dynamic Allocation Mechanics}
	 \subsection{Heap Management Strategies}
	 \section{Multidimensional Mapping}
	 \subsection{Row-Major vs. Column-Major}
	 \subsection{Blocked Memory Layouts}
	 
	 \chapter{Array Indexing Evolution}
	 \section{Address Calculation Mathematics}
	 \subsection{Generalized Dimensional Formula}
	 \section{Bounds Checking Implementations}
	 \subsection{Hardware vs. Software Approaches}
	 \section{Pointer/Array Duality in C}
	 
	 \part{Advanced Array Concepts}
	 \chapter{Low-Level Optimization Techniques}
	 \section{Cache-Aware Array Traversal}
	 \section{SIMD Vectorization Strategies}
	 \section{False Sharing Prevention}
	 
	 \chapter{Theoretical Foundations}
	 \section{Arrays in Automata Theory}
	 \section{Turing Machines with Array Tapes}
	 \section{Chomsky Hierarchy Relationships}
	 
	 \chapter{Specialized Array Architectures}
	 \section{Sparse Array Storage}
	 \subsection{Compressed Sparse Row Format}
	 \section{Jagged Array Implementations}
	 \section{Associative Array Designs}

	 \chapter{Computer Architecture Supplement}
	 \section{From Vacuum Tubes to VLSI}
	 \section{Pipeline Architectures Deep Dive}
	 
	 \chapter{Number System Reference}
	 \section{Positional Number Proofs}
	 \section{Endianness Conversion Algorithms}
	 
	\chapter{Introduction to Arrays}
	\section{Overview}
	\section{Why Use Arrays?}
	\section{History}
	\chapter{Basics of Array Operations}
	\section{Traversal Operation}
	\section{Insertion Operation}
	\section{Deletion Operation}
	\section{Search Operation}
	\section{Sorting Operation}
	\section{Access Operation}
	\chapter{Types and Representations of Arrays}
	\section{Chomsky}
	\section{Types}
	\section{Abstract Arrays}
	\chapter{Memory Layout and Storage}
	\section{Memory Layout of Arrays}
	\section{Memory Segmentation and Bounds Checking}
	\subsection{Memory Segmentation}
	\subsubsection{Hardware Implementation}
	\subsubsection{Segmentation without Paging}
	\subsubsection{Segmentation with Paging}
	\subsubsection{Historical Implementations}
	\subsubsection{x86 Architecture}
	\subsection{Index-Bounds Checking}
	\subsubsection{Range Checking}
	\subsubsection{Index Checking}
	\subsubsection{Hardware Bounds Checking}
	\subsubsection{Support in High-Level Programming Languages}
	\subsubsection{Buffer Overflow}
	\subsubsection{Integer Overflow}
	\chapter{Development of Array Indexing}
	\subsubsection{Address Calculation for Multi-dimensional Arrays}
	\subsubsection{One-Dimensional Array}
	\subsubsection{Two-Dimensional Array}
	\subsubsection{Three-Dimensional Array}
	\subsubsection{Generalizing to a k-Dimensional Array}
	\subsubsection{Examples}
	\chapter{Array Algorithms}
	\section{Sorting Algorithms}
	\section{Searching Algorithms}
	\section{Array Manipulation Algorithms}
	\section{Dynamic Programming and Arrays}
	\chapter{Practical and Advanced Topics}
	\section{Self-Modifying Code in Early Computers}
	\section{Common Array Algorithms}
	\section{Performance Considerations}
	\section{Practical Applications of Arrays}
	\section{Future Trends in Array Handling}
	\chapter{Implementing Arrays in Low-Level Languages}
	\chapter{Static Arrays}
	\section{Single-Dimensional Arrays}
	\subsection{Declaration and Initialization}
	\subsection{Accessing Elements}
	\subsection{Iterating Through an Array}
	\subsection{Common Operations}
	\subsubsection{Insertion}
	\subsubsection{Deletion}
	\subsubsection{Searching}
	\subsection{Memory Considerations}
	\section{Multi-Dimensional Arrays}
	\subsection{2D Arrays}
	\subsubsection{Declaration and Initialization}
	\subsubsection{Accessing Elements}
	\subsubsection{Iterating Through a 2D Array}
	\subsection{3D Arrays and Higher Dimensions}
	\subsubsection{Declaration and Initialization}
	\subsubsection{Accessing Elements}
	\subsubsection{Use Cases and Applications}
	\chapter{Dynamic Arrays}
	\section{Introduction to Dynamic Arrays}
	\subsection{Definition and Overview}
	\subsection{Comparison with Static Arrays}
	
	\section{Single-Dimensional Dynamic Arrays}
	\subsection{Using \texttt{malloc} and \texttt{calloc} in C}
	\subsection{Resizing Arrays with \texttt{realloc}}
	\subsection{Using \texttt{ArrayList} in Java}
	\subsection{Using \texttt{Vector} in C++}
	\subsection{Using \texttt{List} in Python}
	
	\section{Multi-Dimensional Dynamic Arrays}
	\subsection{2D Dynamic Arrays}
	\subsubsection{Creating and Resizing 2D Arrays}
	\subsection{3D and Higher Dimensions}
	\subsubsection{Memory Allocation Techniques}
	\subsubsection{Use Cases and Applications}
	
	\chapter{Advanced Topics in Arrays}
	
	\section{Array Algorithms}
	\subsection{Sorting Algorithms}
	\subsubsection{Bubble Sort}
	\subsubsection{Merge Sort}
	\subsection{Searching Algorithms}
	\subsubsection{Linear Search}
	\subsubsection{Binary Search}
	
	\section{Memory Management in Arrays}
	\subsection{Static vs. Dynamic Memory}
	\subsection{Optimizing Memory Usage}
	
	\section{Handling Large Data Sets}
	\subsection{Efficient Storage Techniques}
	\subsection{Using Arrays in Big Data Applications}
	
	\section{Parallel Processing with Arrays}
	\subsection{Introduction to Parallel Arrays}
	\subsection{Applications in GPU Programming}
	
	\section{Sparse Arrays}
	\subsection{Representation and Usage}
	\subsection{Applications in Data Compression}
	\section{Multidimensional Arrays}
	\section{Jagged Arrays}
	\section{Sparse Arrays}
	\section{Array of Structures vs. Structure of Arrays}
	\section{Array-Based Data Structures}
	
	\chapter{Arrays in Theoretical Computing Paradigms}
	
	\section{Introduction to Theoretical Computing Paradigms}
	\section{Arrays in Turing Machines}
	\section{Arrays in Cellular Automata}
	\section{Arrays in Cellular Automata}
	\section{Arrays in Quantum Computing}
	\section{Arrays in Neural Network Simulations}
	\section{Arrays in Automata Theory}
	\section{Arrays in Hypercomputation Models}
	\section{The Lambda Calculus Perspective on Arrays}
	\section{Arrays in Novel Computational Models}
	
	\chapter{Specialized Arrays and Applications}
	\section{Circular Buffers}
	\section{Circular Arrays}
	\subsection{Implementation and Use Cases}
	\subsection{Applications in Buffer Management}
	
	\section{Dynamic Buffering and Arrays}
	\subsection{Dynamic Circular Buffers}
	\subsection{Handling Streaming Data}
	
	\section{Jagged Arrays}
	\subsection{Definition and Usage}
	\subsection{Applications in Database Management}
	
	\section{Bit Arrays (Bitsets)}
	\subsection{Introduction and Representation}
	\subsection{Applications in Cryptography}
	\section{Circular Buffers}
	\section{Priority Queues}
	\section{Hash Tables}
	\section{Bloom Filters}
	\section{Bit Arrays and Bit Vectors}
	
	\chapter{Linked Lists}
	\section{Overview}
	\section{Singly Linked Lists}
	\section{Doubly Linked Lists}
	\section{Circular Linked Lists}
	\section{Comparison with Arrays}
	
	\chapter{Array-Based Algorithms}
	\section{Sorting Algorithms}
	\section{Searching Algorithms}
	\section{Array Manipulation Algorithms}
	\section{Dynamic Programming and Arrays}
	
	\chapter{Performance Analysis}
	\section{Time Complexity of Array Operations}
	\section{Space Complexity Considerations}
	\section{Cache Performance and Optimization}
	
	\chapter{Memory Management}
	\section{Memory Allocation Strategies}
	\section{Garbage Collection}
	\section{Manual Memory Management in Low-Level Languages}
	
	\chapter{Error Handling and Debugging}
	\section{Common Errors with Arrays}
	\section{Bounds Checking Techniques}
	\section{Debugging Tools and Strategies}
	
	\chapter{Optimization Techniques for Arrays}
	\section{Optimizing Array Traversal}
	\section{Minimizing Cache Misses}
	\section{Loop Unrolling}
	\section{Vectorization}
	\section{Memory Access Patterns}
	\section{Reducing Memory Fragmentation}
	
	\chapter{Concurrency and Parallelism}
	\section{Concurrent Array Access}
	\section{Parallel Array Processing}
	\section{Synchronization Techniques}
	
	\chapter{Applications in Modern Software Development}
	\section{Arrays in Graphics and Game Development}
	\section{Arrays in Scientific Computing}
	\section{Arrays in Data Analysis and Machine Learning}
	\section{Arrays in Embedded Systems}
	
	\chapter{Arrays in High-Performance Computing (HPC)}
	\section{Introduction to HPC Arrays}
	\section{Distributed Arrays}
	\section{Parallel Processing with Arrays}
	\section{Arrays in GPU Computing}
	\section{Multi-threaded Array Operations}
	\section{Handling Arrays in Cloud Computing}
	
	\chapter{Arrays in Functional Programming}
	\section{Immutable Arrays}
	\section{Persistent Arrays}
	\section{Arrays in Functional Languages (Haskell, Erlang, etc.)}
	\section{Functional Array Operations}
	
	\chapter{Arrays in Machine Learning and Data Science}
	\section{Numerical Arrays}
	\section{Handling Large Datasets with Arrays}
	\section{Arrays in Tensor Operations}
	\section{Arrays in Dataframes}
	\section{Optimization of Array-Based Algorithms in ML}
	
	\chapter{Advanced Memory Management in Arrays}
	\section{Memory Pools}
	\section{Dynamic Memory Allocation Strategies}
	
	\chapter{Data Structures Derived from Arrays}
	\section{Stacks}
	\section{Queues}
	\section{Heaps}
	\section{Hash Tables}
	\section{Trees Implemented Using Arrays}
	\section{Graphs Implemented Using Arrays}
	\section{Dynamic Arrays as Building Blocks}
	
	\chapter{Best Practices and Common Pitfalls in Array Usage}
	\section{Avoiding Out-of-Bounds Errors}
	\section{Efficient Initialization}
	\section{Choosing the Right Array Type}
	\section{Debugging and Testing Arrays}
	\section{Avoiding Memory Leaks}
	\section{Ensuring Portability Across Platforms}
	
	\chapter{Historical Perspectives and Evolution}
	\section{Custom Memory Allocators}\section{Early Implementations}
	\section{Array Storage on Disk}\section{Evolution of Array Data Structures}
	\section{Impact on Programming Languages and Paradigms}
	
	\chapter{Future Trends in Array Handling}
	\section{Emerging Data Structures}
	\section{Quantum Computing and Arrays}
	\section{Bioinformatics Applications}
	\section{Big Data and Arrays}
	\section{Arrays in Emerging Programming Paradigms}
	\chapter{Appendices}
	\section{Glossary of Terms}
	\section{Bibliography}
	\section{Index}
	
% References
\printbibliography[heading=bibintoc]

% Index
\printindex
\end{document}